% Môn: Vật lý
% Lớp: 10
% Chương: Chương I. Mở đầu
% Bài: L10C1 Bài 2. Các quy tắc an toàn trong phòng thực hành Vật lí · An toàn khi sử dụng thiết bị điện
% Số câu: 185
% App đọc file này trực tiếp — sửa rồi Commit trên GitHub, bấm Đồng bộ GitHub .tex

% L10C1 Bài 2. Các quy tắc an toàn trong phòng thực hành Vật lí



% ===== Câu 12 =====
% ID: L10C1B2-06-DS
% Mức: NB
\dangbt{An toàn khi sử dụng thiết bị điện}
\begin{ex}
% ID: L10C1B2-06-DS
% Mức: NB
Trong tình huống khi phát hiện thiết bị có dấu hiệu bất thường đối với thiết bị điện, hãy đánh giá các phát biểu sau.
\choiceTF
{\True Cần ngắt nguồn, kiểm tra dây dẫn và dùng tay khô.}
{Có thể chạm vào phần dẫn điện khi mạch đang hoạt động nếu thao tác diễn ra nhanh.}
{\True Phải báo cho giáo viên hoặc người phụ trách khi phát hiện nguy cơ vượt khả năng kiểm soát.}
{\True Chỉ tiếp tục thí nghiệm sau khi nguyên nhân nguy hiểm đã được loại bỏ hoặc kiểm soát.}
\loigiai{
\begin{itemchoice}
\item (Đúng) Cần ngắt nguồn, kiểm tra dây dẫn và dùng tay khô. (Vì): vì đây là biện pháp an toàn của dạng thiết bị điện; b sai vì hành vi “chạm vào phần dẫn điện khi mạch đang hoạt động” vẫn nguy hiểm; c và d đúng do sự cố phải được báo cáo, cô lập và kiểm soát trước khi tiếp tục
\item (Sai) Có thể chạm vào phần dẫn điện khi mạch đang hoạt động nếu thao tác diễn ra nhanh.
\item (Đúng) Phải báo cho giáo viên hoặc người phụ trách khi phát hiện nguy cơ vượt khả năng kiểm soát.
\item (Đúng) Chỉ tiếp tục thí nghiệm sau khi nguyên nhân nguy hiểm đã được loại bỏ hoặc kiểm soát.
\end{itemchoice}
}
\end{ex}

% ===== Câu 13 =====
% ID: L10C1B2-07-DS
% Mức: NB
\begin{ex}
% ID: L10C1B2-07-DS
% Mức: NB
Trong tình huống trong lúc thay đổi cách bố trí dụng cụ đối với thiết bị điện, hãy đánh giá các phát biểu sau.
\choiceTF
{\True Cần ngắt nguồn, kiểm tra dây dẫn và dùng tay khô.}
{Có thể chạm vào phần dẫn điện khi mạch đang hoạt động nếu thao tác diễn ra nhanh.}
{\True Phải báo cho giáo viên hoặc người phụ trách khi phát hiện nguy cơ vượt khả năng kiểm soát.}
{\True Chỉ tiếp tục thí nghiệm sau khi nguyên nhân nguy hiểm đã được loại bỏ hoặc kiểm soát.}
\loigiai{
\begin{itemchoice}
\item (Đúng) Cần ngắt nguồn, kiểm tra dây dẫn và dùng tay khô. (Vì): vì đây là biện pháp an toàn của dạng thiết bị điện; b sai vì hành vi “chạm vào phần dẫn điện khi mạch đang hoạt động” vẫn nguy hiểm; c và d đúng do sự cố phải được báo cáo, cô lập và kiểm soát trước khi tiếp tục
\item (Sai) Có thể chạm vào phần dẫn điện khi mạch đang hoạt động nếu thao tác diễn ra nhanh.
\item (Đúng) Phải báo cho giáo viên hoặc người phụ trách khi phát hiện nguy cơ vượt khả năng kiểm soát.
\item (Đúng) Chỉ tiếp tục thí nghiệm sau khi nguyên nhân nguy hiểm đã được loại bỏ hoặc kiểm soát.
\end{itemchoice}
}
\end{ex}

% ===== Câu 14 =====
% ID: L10C1B2-07-TL
% Mức: TH
\begin{bt}
% ID: L10C1B2-07-TL
% Mức: TH
Trình bày nguyên tắc cốt lõi khi làm việc với thiết bị điện và giải thích vì sao phải tuân thủ.
\loigiai{
Nguyên tắc cốt lõi là ngắt nguồn, kiểm tra dây dẫn và dùng tay khô. Phải tuân thủ vì tránh điện giật, chập mạch và hư hỏng thiết bị; đồng thời luôn dừng và báo người phụ trách khi điều kiện không bảo đảm.
}
\end{bt}

% ===== Câu 15 =====
% ID: L10C1B2-07-TLN
% Mức: NB
\begin{ex}
% ID: L10C1B2-07-TLN
% Mức: NB
Một quy trình về thiết bị điện gồm: (1) nhận diện nguy cơ; (2) ngắt nguồn, kiểm tra dây dẫn và dùng tay khô; (3) chạm vào phần dẫn điện khi mạch đang hoạt động; (4) báo giáo viên khi có bất thường; (5) chỉ tiếp tục khi nguy cơ đã được kiểm soát. Có bao nhiêu bước phù hợp quy tắc an toàn? Chỉ ghi số.
\shortans{4}
\loigiai{
Các bước (1), (2), (4), (5) phù hợp; bước (3) không an toàn. Có 4 bước phù hợp.
}
\end{ex}

% ===== Câu 16 =====
% ID: L10C1B2-08-TL
% Mức: TH
\begin{bt}
% ID: L10C1B2-08-TL
% Mức: TH
Phân tích hành vi “chạm vào phần dẫn điện khi mạch đang hoạt động” và đề xuất cách thay thế an toàn.
\loigiai{
Hành vi “chạm vào phần dẫn điện khi mạch đang hoạt động” có thể gây tai nạn vì làm mất kiểm soát nguồn nguy hiểm. Cần dừng thao tác, cô lập nguy cơ và thay bằng biện pháp: ngắt nguồn, kiểm tra dây dẫn và dùng tay khô.
}
\end{bt}

% ===== Câu 18 =====
% ID: L10C1B2-09-DS
% Mức: TH
\begin{ex}
% ID: L10C1B2-09-DS
% Mức: TH
Trong tình huống khi một bạn chưa đọc hướng dẫn đối với thiết bị điện, hãy đánh giá các phát biểu sau.
\choiceTF
{\True Cần ngắt nguồn, kiểm tra dây dẫn và dùng tay khô.}
{Có thể chạm vào phần dẫn điện khi mạch đang hoạt động nếu thao tác diễn ra nhanh.}
{\True Phải báo cho giáo viên hoặc người phụ trách khi phát hiện nguy cơ vượt khả năng kiểm soát.}
{\True Chỉ tiếp tục thí nghiệm sau khi nguyên nhân nguy hiểm đã được loại bỏ hoặc kiểm soát.}
\loigiai{
\begin{itemchoice}
\item (Đúng) Cần ngắt nguồn, kiểm tra dây dẫn và dùng tay khô. (Vì): vì đây là biện pháp an toàn của dạng thiết bị điện; b sai vì hành vi “chạm vào phần dẫn điện khi mạch đang hoạt động” vẫn nguy hiểm; c và d đúng do sự cố phải được báo cáo, cô lập và kiểm soát trước khi tiếp tục
\item (Sai) Có thể chạm vào phần dẫn điện khi mạch đang hoạt động nếu thao tác diễn ra nhanh.
\item (Đúng) Phải báo cho giáo viên hoặc người phụ trách khi phát hiện nguy cơ vượt khả năng kiểm soát.
\item (Đúng) Chỉ tiếp tục thí nghiệm sau khi nguyên nhân nguy hiểm đã được loại bỏ hoặc kiểm soát.
\end{itemchoice}
}
\end{ex}

% ===== Câu 20 =====
% ID: L10C1B2-10-DS
% Mức: VD
\begin{ex}
% ID: L10C1B2-10-DS
% Mức: VD
Trong tình huống khi khu vực làm việc bị ướt đối với thiết bị điện, hãy đánh giá các phát biểu sau.
\choiceTF
{\True Cần ngắt nguồn, kiểm tra dây dẫn và dùng tay khô.}
{Có thể chạm vào phần dẫn điện khi mạch đang hoạt động nếu thao tác diễn ra nhanh.}
{\True Phải báo cho giáo viên hoặc người phụ trách khi phát hiện nguy cơ vượt khả năng kiểm soát.}
{\True Chỉ tiếp tục thí nghiệm sau khi nguyên nhân nguy hiểm đã được loại bỏ hoặc kiểm soát.}
\loigiai{
\begin{itemchoice}
\item (Đúng) Cần ngắt nguồn, kiểm tra dây dẫn và dùng tay khô. (Vì): vì đây là biện pháp an toàn của dạng thiết bị điện; b sai vì hành vi “chạm vào phần dẫn điện khi mạch đang hoạt động” vẫn nguy hiểm; c và d đúng do sự cố phải được báo cáo, cô lập và kiểm soát trước khi tiếp tục
\item (Sai) Có thể chạm vào phần dẫn điện khi mạch đang hoạt động nếu thao tác diễn ra nhanh.
\item (Đúng) Phải báo cho giáo viên hoặc người phụ trách khi phát hiện nguy cơ vượt khả năng kiểm soát.
\item (Đúng) Chỉ tiếp tục thí nghiệm sau khi nguyên nhân nguy hiểm đã được loại bỏ hoặc kiểm soát.
\end{itemchoice}
}
\end{ex}

% ===== Câu 29 =====
% ID: L10C1B2-11-TL
% Mức: VD
\begin{bt}
% ID: L10C1B2-11-TL
% Mức: VD
So sánh biện pháp phòng ngừa và biện pháp ứng phó đối với nguy cơ thuộc nhóm thiết bị điện.
\loigiai{
Phòng ngừa diễn ra trước sự cố, tập trung vào kiểm tra, loại bỏ hoặc giảm nguy cơ. Ứng phó diễn ra khi sự cố đã hoặc sắp xảy ra, gồm dừng hoạt động, cô lập nguồn nguy hiểm, báo cáo và sơ cứu. Hai nhóm biện pháp bổ sung cho nhau.
}
\end{bt}

% ===== Câu 30 =====
% ID: L10C1B2-11-TN
% Mức: NB
\begin{ex}
% ID: L10C1B2-11-TN
% Mức: NB
Trước khi bắt đầu thí nghiệm liên quan đến thiết bị điện, hành động nào an toàn nhất?
\choice
{\True ngắt nguồn, kiểm tra dây dẫn và dùng tay khô}
{chạm vào phần dẫn điện khi mạch đang hoạt động}
{tiếp tục thao tác rồi báo cáo sau}
{nhờ bạn khác làm thay mà không kiểm tra điều kiện an toàn}
\loigiai{
Hành động đúng là “ngắt nguồn, kiểm tra dây dẫn và dùng tay khô” vì tránh điện giật, chập mạch và hư hỏng thiết bị. Chọn A.
}
\end{ex}

% ===== Câu 40 =====
% ID: L10C1B2-12-TL
% Mức: VDC
\begin{bt}
% ID: L10C1B2-12-TL
% Mức: VDC
Đề xuất cách tổ chức nhóm học sinh để thực hiện nhiệm vụ về thiết bị điện an toàn và có thể kiểm tra được.
\loigiai{
Phân công trưởng nhóm kiểm tra điều kiện, người thao tác, người quan sát an toàn và người ghi chép. Dùng bảng kiểm, xác nhận trước mỗi bước, quy định tín hiệu dừng và báo ngay bất thường.
}
\end{bt}

% ===== Câu 41 =====
% ID: L10C1B2-12-TN
% Mức: NB
\begin{ex}
% ID: L10C1B2-12-TN
% Mức: NB
Khi phát hiện thiết bị có dấu hiệu bất thường liên quan đến thiết bị điện, hành động nào an toàn nhất?
\choice
{chạm vào phần dẫn điện khi mạch đang hoạt động}
{tiếp tục thao tác rồi báo cáo sau}
{nhờ bạn khác làm thay mà không kiểm tra điều kiện an toàn}
{\True ngắt nguồn, kiểm tra dây dẫn và dùng tay khô}
\loigiai{
Hành động đúng là “ngắt nguồn, kiểm tra dây dẫn và dùng tay khô” vì tránh điện giật, chập mạch và hư hỏng thiết bị. Chọn D.
}
\end{ex}

% ===== Câu 53 =====
% ID: L10C1B2-13-TN
% Mức: NB
\begin{ex}
% ID: L10C1B2-13-TN
% Mức: NB
Trong lúc thay đổi cách bố trí dụng cụ liên quan đến thiết bị điện, hành động nào an toàn nhất?
\choice
{tiếp tục thao tác rồi báo cáo sau}
{nhờ bạn khác làm thay mà không kiểm tra điều kiện an toàn}
{\True ngắt nguồn, kiểm tra dây dẫn và dùng tay khô}
{chạm vào phần dẫn điện khi mạch đang hoạt động}
\loigiai{
Hành động đúng là “ngắt nguồn, kiểm tra dây dẫn và dùng tay khô” vì tránh điện giật, chập mạch và hư hỏng thiết bị. Chọn C.
}
\end{ex}

% ===== Câu 65 =====
% ID: L10C1B2-14-TN
% Mức: TH
\begin{ex}
% ID: L10C1B2-14-TN
% Mức: TH
Sau khi hoàn thành phép đo liên quan đến thiết bị điện, hành động nào an toàn nhất?
\choice
{nhờ bạn khác làm thay mà không kiểm tra điều kiện an toàn}
{\True ngắt nguồn, kiểm tra dây dẫn và dùng tay khô}
{chạm vào phần dẫn điện khi mạch đang hoạt động}
{tiếp tục thao tác rồi báo cáo sau}
\loigiai{
Hành động đúng là “ngắt nguồn, kiểm tra dây dẫn và dùng tay khô” vì tránh điện giật, chập mạch và hư hỏng thiết bị. Chọn B.
}
\end{ex}

% ===== Câu 77 =====
% ID: L10C1B2-16-TN
% Mức: TH
\begin{ex}
% ID: L10C1B2-16-TN
% Mức: TH
Khi khu vực làm việc bị ướt liên quan đến thiết bị điện, hành động nào an toàn nhất?
\choice
{chạm vào phần dẫn điện khi mạch đang hoạt động}
{tiếp tục thao tác rồi báo cáo sau}
{nhờ bạn khác làm thay mà không kiểm tra điều kiện an toàn}
{\True ngắt nguồn, kiểm tra dây dẫn và dùng tay khô}
\loigiai{
Hành động đúng là “ngắt nguồn, kiểm tra dây dẫn và dùng tay khô” vì tránh điện giật, chập mạch và hư hỏng thiết bị. Chọn D.
}
\end{ex}

% ===== Câu 80 =====
% ID: L10C1B2-17-TN
% Mức: TH
\begin{ex}
% ID: L10C1B2-17-TN
% Mức: TH
Khi cần rời bàn thí nghiệm liên quan đến thiết bị điện, hành động nào an toàn nhất?
\choice
{tiếp tục thao tác rồi báo cáo sau}
{nhờ bạn khác làm thay mà không kiểm tra điều kiện an toàn}
{\True ngắt nguồn, kiểm tra dây dẫn và dùng tay khô}
{chạm vào phần dẫn điện khi mạch đang hoạt động}
\loigiai{
Hành động đúng là “ngắt nguồn, kiểm tra dây dẫn và dùng tay khô” vì tránh điện giật, chập mạch và hư hỏng thiết bị. Chọn C.
}
\end{ex}

% ===== Câu 83 =====
% ID: L10C1B2-18-TN
% Mức: VD
\begin{ex}
% ID: L10C1B2-18-TN
% Mức: VD
Khi có người không đeo phương tiện bảo hộ liên quan đến thiết bị điện, hành động nào an toàn nhất?
\choice
{nhờ bạn khác làm thay mà không kiểm tra điều kiện an toàn}
{\True ngắt nguồn, kiểm tra dây dẫn và dùng tay khô}
{chạm vào phần dẫn điện khi mạch đang hoạt động}
{tiếp tục thao tác rồi báo cáo sau}
\loigiai{
Hành động đúng là “ngắt nguồn, kiểm tra dây dẫn và dùng tay khô” vì tránh điện giật, chập mạch và hư hỏng thiết bị. Chọn B.
}
\end{ex}

% ===== Câu 85 =====
% ID: L10C1B2-19-TN
% Mức: VD
\begin{ex}
% ID: L10C1B2-19-TN
% Mức: VD
Khi kết quả đo khác dự kiến liên quan đến thiết bị điện, hành động nào an toàn nhất?
\choice
{\True ngắt nguồn, kiểm tra dây dẫn và dùng tay khô}
{chạm vào phần dẫn điện khi mạch đang hoạt động}
{tiếp tục thao tác rồi báo cáo sau}
{nhờ bạn khác làm thay mà không kiểm tra điều kiện an toàn}
\loigiai{
Hành động đúng là “ngắt nguồn, kiểm tra dây dẫn và dùng tay khô” vì tránh điện giật, chập mạch và hư hỏng thiết bị. Chọn A.
}
\end{ex}

% ===== Câu 87 =====
% ID: L10C1B2-20-TN
% Mức: VD
\begin{ex}
% ID: L10C1B2-20-TN
% Mức: VD
Khi giáo viên yêu cầu dừng thao tác liên quan đến thiết bị điện, hành động nào an toàn nhất?
\choice
{chạm vào phần dẫn điện khi mạch đang hoạt động}
{tiếp tục thao tác rồi báo cáo sau}
{nhờ bạn khác làm thay mà không kiểm tra điều kiện an toàn}
{\True ngắt nguồn, kiểm tra dây dẫn và dùng tay khô}
\loigiai{
Hành động đúng là “ngắt nguồn, kiểm tra dây dẫn và dùng tay khô” vì tránh điện giật, chập mạch và hư hỏng thiết bị. Chọn D.
}
\end{ex}

